\chapter{はじめに}\label{cha:Introduction}

ソフトウェア開発において、システム仕様書は開発の指針となる極めて重要なドキュメントである。
特に、顧客の曖昧な要望を具体的な実装レベルの仕様に落とし込む過程において、仕様書の品質はプロジェクトの成否を左右する。

システム仕様書の作成は、大きく「章立て（構造）の検討」と「各章の執筆」の2段階に分けられる。
このうち、最初のステップである「章立ての検討」は、必要な要件を網羅し、論理的な構成を決定する最も重要な工程である。
もしこの段階で重要な項目の欠落や順序の不整合が生じると、後続の執筆作業すべてに悪影響を及ぼし、多大な手戻りが発生する。
しかし、この工程は高度な知識と経験を要するため、特に経験の浅いエンジニアにとっては難易度が高い作業となっている。

近年、大規模言語モデル（LLM）を用いたドキュメント作成の自動化が注目されているが、長文の仕様書を一度に生成させようとすると、文脈の維持が難しく、構造が破綻したりハルシネーション（事実に基づかない生成）が発生したりする問題がある。
そのため、実用的な品質の仕様書を得るためには、全文を一括生成するのではなく、段階的なアプローチが必要となる。

そこで本研究では、システム仕様書作成の第一歩であり、品質の根幹となる「章立て（構造）の生成」に着目する。
プロンプトエンジニアリングの技術を適用し、システム概要から網羅的かつ論理的な仕様書の章立てを自動生成する手法を提案する。
これにより、仕様書作成における構造的な欠陥を早期に防ぎ、ドキュメント品質の均質化と作成効率の向上を図ることを目的とする。

本論文の構成は以下の通りである。
第\ref{cha:Preparation}章では、本研究の背景知識としてプロンプトエンジニアリングの技術やシステム仕様書の定義について述べる。
第\ref{cha:Proposal}章では、提案するプロンプトの構成とその利用プロセスについて詳述する。
第\ref{cha:Application}章では、提案手法を実際の開発事例に適用した結果を示す。
第\ref{cha:Evaluation}章では、生成された章立ての品質評価と比較実験の結果について述べる。
第\ref{cha:Discussion}章では、得られた結果の考察と今後の課題について論じる。
最後に第\ref{cha:Conclusion}章で本論文をまとめる。